
\FloatBarrier
\section{SEM for building material characterisation}
% imaging overview paper \cite{qoku.2023}
The fist electron microscope was a \ac{TEM}, constructed by \textsc{Ernst Ruska} and \textsc{Max Knoll} in 1931 \cite{knoll.1932} which is based on the transmission of electrons trough the specimen.
It was first applied in 1939 by \textcite{rapczewski.1939} to image hydration products of \ac{C3S}.
In the 1950s and 1960s, electron microscopy in combination with electron diffraction developed to become a useful and reliable tool for the analysis of cement and its hydration products \cite{grudemo.1960, grudemo.1965}.
However, it had multiple drawbacks as the requirement of very small (powdered) and thin specimens to be able to penetrate the specimen with the electron beam.

In 1937, \textsc{Manfred von Ardenne} build the first \ac{SEM} which enabled to also the surface of specimens \cite{vonardenne.1938}, allowing for a wider range of preparation methods.
It took until 1965, when the first commercial \ac{SEM} made by \textsc{Cambridge Instrument Co.} was available.
Only a year later, \textcite{chatterji.1966} used this \ac{SEM} to observe fractured specimens of hydrated cement.
They already observed the formation of the needle-like structure of \ac{CSH}.
Today, \ac{SEM} is a widely used tool for the analysis of building materials.
Modern \ac{SEM}s are equipped with various detectors, like \ac{SE}, \ac{BSE}, \ac{EDX} and \ac{EBSD} detectors.
They allow to observe the chemical composition, the crystallographic orientation and the topography of the specimen.

In 1974/75 the \ac{FIB} technique was introduced by \textcite{levi-setti.1974} and \textcite{orloff.1975} for microprobing applications.
\ac{FIB} systems are commercially available since the 1990s and today combined \ac{FIB}-\ac{SEM} devices are available.
The \ac{FIB} technique is widely used for the preparation of samples for \ac{SEM} and \ac{TEM}.
In contrast to electron microscopy, \ac{FIB} uses a beam of ions instead of electrons which interact with the specimen surface.
While it is possible to obtain an image of the specimen surface, the main application is the preparation of samples for \ac{SEM} and \ac{TEM} due to the abrading effect of the ions.
\cite{holzer.2006, munch.2006} analyzed particular systems using \ac{FIB} serial sectioning.
This technique allows to obtain a digital 3D volume of the specimen.
They also established a segmentation workflow proposing to analyse cryo-stabilized cement suspensions which they published shortly after \textcite{holzer.2007}.

\textcite{holzer.2006a} used \ac{FIBnt} very early to cut a large volume (typically \num{46}\,$\times$\,\num{32}\,$\times$\,\SI{25}{\cubic\micro\meter}) of a cement paste sample with a resolution of \SI{74}{\nano\meter}.

-cryoSEM+nT? -> \cite{holzer.2007,desbois.2008}



%plasma fib: \cite{dong.2022} (larger volume than Ga-FIB)
% FIB-SEM of OPC: \cite{jiang.2023}
% \cite{song.2019}
% \cite{holzer.2006a} very early study of a very large volume (46*32*25μm) @voxel size of 74nm - good quality
% \cite{mac.2021} very good paper describing the processing of FIB datasets + BIB-tomography
% \cite{munch.2008} FIB vs MIP!!! (11.9 x10.7 x5.6 µm)
% shitty dataset \cite{yang.2021}

% \cite{Lavrov.2019} - ct of cement paste - ~ 20µm ITZ


% was gibts für abbildungsmodi - praktikumsanleitungen - Karen Microstructure analysis
%state of the art Bildanalyse - stutzman, fabien, Holzer, Münch, Herr Ulm/MITv -> Segmentierung/Kanten




% \begin{figure}[htbp]
% 	\centerline{\includegraphics[width=.49\textwidth]{hydrate-rim/hadley_grains_detail.pdf}}
% 	\caption{
% 	\label{fig:hadley2}}
% \end{figure}